% La griglia studiata in Sudoku.dfw e i determinanti dei nove blocchi.
\begin{tikzpicture}[scale=0.56, every node/.style={font=\small}]

  % ---- la griglia ---------------------------------------------------
  \begin{scope}
    \def\griglia{
      {9,1,3,7,8,6,5,4,2},
      {4,8,7,5,2,3,1,9,6},
      {2,6,5,4,1,9,3,8,7},
      {8,5,6,3,7,4,2,1,9},
      {7,9,2,1,6,5,4,3,8},
      {3,4,1,2,9,8,7,6,5},
      {6,3,4,9,5,7,8,2,1},
      {5,2,9,8,4,1,6,7,3},
      {1,7,8,6,3,2,9,5,4}}

    \foreach \bi in {0,1,2}{
      \foreach \bj in {0,1,2}{
        \pgfmathtruncatemacro{\tinta}{mod(\bi+\bj,2)}
        \ifnum\tinta=0
          \fill[grigiocarta] (3*\bj,-3*\bi) rectangle ++(3,-3);
        \fi
      }}

    \foreach \r [count=\i from 0] in \griglia {
      \foreach \v [count=\j from 0] in \r {
        \node[inchiostro] at (\j+0.5,-\i-0.5) {\v};
      }}

    \foreach \k in {0,...,9}{
      \draw[grigiotratto,line width=0.5pt] (\k,0) -- (\k,-9);
      \draw[grigiotratto,line width=0.5pt] (0,-\k) -- (9,-\k);
    }
    \foreach \k in {0,3,6,9}{
      \draw[blunotte,line width=1.1pt] (\k,0) -- (\k,-9);
      \draw[blunotte,line width=1.1pt] (0,-\k) -- (9,-\k);
    }
    \node[grigiotesto,font=\footnotesize\sffamily,anchor=north] at (4.5,-9.3)
      {la griglia};
  \end{scope}

  % ---- i determinanti dei blocchi -----------------------------------
  \begin{scope}[xshift=11.2cm]
    % celle larghe 2 e alte 2.5, per allinearsi ai blocchi della griglia
    \foreach \d [count=\n from 0] in {0,-177,81,9,11,-3,-243,-5,77}{
      \pgfmathtruncatemacro{\bi}{div(\n,3)}
      \pgfmathtruncatemacro{\bj}{mod(\n,3)}
      \node[rossoterra] at (2*\bj+1,{-0.75-2.5*\bi}) {$\d$};
    }
    \foreach \k in {2,4}{
      \draw[grigiotratto,line width=0.5pt] (\k,0.5) -- (\k,-7);
    }
    \foreach \k in {1,2}{
      \draw[grigiotratto,line width=0.5pt] (0,{0.5-2.5*\k}) -- (6,{0.5-2.5*\k});
    }
    \draw[blunotte,line width=1.1pt] (0,0.5) rectangle (6,-7);
    \node[grigiotesto,font=\footnotesize\sffamily,anchor=north,align=center]
      at (3,-7.4)
      {determinante di ciascun blocco\\[2pt] somma: $-250$};
  \end{scope}

  % ---- didascalia complessiva ---------------------------------------
  \node[grigiotesto,font=\footnotesize\sffamily,anchor=north] at (8.6,-10.6)
    {determinante della griglia intera: $-2\,975\,130=-2\cdot3^4\cdot5\cdot3673$};
\end{tikzpicture}
